\documentclass[a4paper, 12pt]{article}
\usepackage[T2A]{fontenc}
\usepackage[english,russian]{babel}
\usepackage{amsmath}
\usepackage{amsthm}
\usepackage{amsfonts}
\usepackage{amssymb}
\usepackage{pgfplots}

\usepackage{enumitem}
% Абзацы и геометрия страницы
\usepackage{indentfirst}
%\usepackage{geometry}
\usepackage[
    top=2cm,
    bottom=2cm,
    left=1cm,
    right=1cm
]{geometry}
\begin{document}


\begin{center}
    МИНИСТЕРСТВО НАУКИ И ВЫСШЕГО ОБРАЗОВАНИЯ РОССИЙСКОЙ ФЕДЕРАЦИИ

    ГОСУДАРСТВЕННОЕ ОБРАЗОВАТЕЛЬНОЕ УЧРЕЖДЕНИЕ ВЫСШЕГО ПРОФЕССИОНАЛЬНОГО ОБРАЗОВАНИЯ

    РЯЗАНСКИЙ ГОСУДАРСТВЕННЫЙ РАДИОТЕХНИЧЕСКИЙ УНИВЕРСИТЕТ

    ИМ. В. Ф. УТКИНА

    КАФЕДРА ВЫЧИСЛИТЕЛЬНОЙ И ПРИКЛАДНОЙ МАТЕМАТИКИ
\end{center}

\vspace{5 cm}

\begin{center}
    ЛАБОРАТОРНАЯ РАБОТА № 4

    ПО ДИСЦИПЛИНЕ "ТЕОРЕТИЧЕСКАЯ ИНФОРМАТИКА"

    НА ТЕМУ "ПРЕДСТАВЛЕНИЕ ЦЕЛЫХ ЧИСЕЛ В ЭВМ. ЧАСТЬ 2"

    ВАРИАНТ 14

\end{center}

\vspace{7 cm}

\begin{flushright}
    Выполнил:

    Студент группы 6413

    Косяков Д.А.

    Проверил

    Доцент кафедры ВПМ

    Никичкин Борис Владимирович
\end{flushright}

\begin{center}
    РЯЗАНЬ 2026
\end{center}

\pagebreak

\textbf{ Цель работы:} посчитать сумму целых чисел, представленных в дополнительном коде в шестнадцатеричной системе счисления.

\textbf{ Задание:} перевести числа А и В в десятичную систему счисления,
сложить их и результат записать в двоичной системе счисления.
Числа А и В представлены в дополнительном коде.

\textbf{ Выполнение задания:} заданы $A_{16}$ = 51 и $B_{16}$ = 90.

Переведем число А из шестнадцатеричной системы счисления в десятичную, минуя двоичную.
Старшая цифра 5 меньше 8, значит знаковый бит равен 0 --- число положительно:
$$A_{10} = 5 \cdot 16^1 + 1 \cdot 16^0 = 80 + 1 = 81$$

Переведем число В из шестнадцатеричной системы счисления в десятичную, минуя двоичную.
Старшая цифра 9 больше либо равна 8, значит знаковый бит равен 1 --- число отрицательно.
Сначала переведем код числа В в десятичную систему без учета знака:
$$90_{16} = 9 \cdot 16^1 + 0 \cdot 16^0 = 144$$

Число восьмиразрядное, а в дополнительном коде отрицательное число $x$ хранится как $2^8 + x = 256 + x$.
Значит, чтобы получить само число, нужно из кода вычесть $2^8 = 256$:
$$B_{10} = 144 - 256 = -112$$

Проверим модуль, оставаясь в шестнадцатеричной системе. Инверсии всех битов соответствует замена каждой
шестнадцатеричной цифры $d$ на $F - d$: $F - 9 = 6$, $F - 0 = F$, получаем $6F_{16}$. Прибавляем 1:
$$6F_{16} + 1 = 70_{16} = 7 \cdot 16^1 + 0 \cdot 16^0 = 112$$
$$|B_{10}| = 112, \quad B_{10} = -112$$

Вычисляем сумму чисел $A_{10}$ и $B_{10}$:
$$A_{10} + B_{10} = 81 + (-112) = -31$$

Переведем модуль суммы в двоичную систему счисления делением на 2:

$$31 : 2 = 15, \text{ остаток } 1$$
$$15 : 2 = 7, \text{ остаток } 1$$
$$7 : 2 = 3, \text{ остаток } 1$$
$$3 : 2 = 1, \text{ остаток } 1$$
$$1 : 2 = 0, \text{ остаток } 1$$

Если $31_{10} = 11111_2$, значит $-31_{10} = -11111_2$.

В восьмиразрядном дополнительном коде: $31_{10} = 00011111_2$, инверсия --- 11100000, инверсия + 1 = 11100001.
Значит $-31_{10} = 11100001_2$ (доп. код).

Проверка: $51_{16} + 90_{16} = E1_{16} = 11100001_2$.
\vspace{1cm}

\textbf{Отчет:}
\begin{enumerate}[nosep]
    \item А = 51, В = 90.
    \item $A_{10} = 81$, $B_{10} = -112$.
    \item $A_{10} + B_{10} = -31$.
    \item $-31_{10} = -11111_2$, в дополнительном коде 11100001.
\end{enumerate}

\end{document}
